\chapter{Quando buscar e quando decidir: acrescentamos o \nt{NORA}}
\chaptermark{Acrescentamos o \nt{NORA}}
\label{sec:nora}

\section{O que falta}

No capítulo~\ref{sec:dofa}, a rede aprendia graças ao ruído: os desvios aleatórios de atividade eram a busca, e o \nt{DOPA} informava se eles tinham melhorado as coisas. Mas lá ficou sem definir um parâmetro: \emph{quanto} buscar. Se o ruído é forte demais, a rede se debate e não aproveita o que já sabe. Se é fraco demais, ela escolhe sempre a mesma coisa e não percebe que, em algum lugar perto, as coisas melhoraram. No aprendizado por reforço, isso se chama o dilema entre aproveitamento e busca, e todo algoritmo tem um botão para ele. No capítulo~\ref{sec:neurontypes}, supusemos que no cérebro quem gira esse botão é o \nt{NORA}.

Os neurônios noradrenérgicos no ser humano são algumas dezenas de milhares (tabela~\ref{tab:types}), quase todos reunidos num único grupo pequeno, e suas saídas se espalham por praticamente o cérebro inteiro, embora partes diferentes desse grupo, como se descobriu, atendam em parte regiões diferentes~\cite{Poe2020}. É um canal de difusão ainda mais estreito que o \nt{DOPA}. Ele funciona em dois modos~\cite{AstonJones2005}. No modo \emph{de fundo}, os neurônios emitem disparos regulares, com taxa de frações a alguns disparos por segundo, e essa taxa muda lentamente com a vigília. No modo \emph{de rajadas}, eles respondem com uma rajada curta a um evento importante para a tarefa atual, mais ou menos no instante da decisão. Um ponto de teste cômodo, mas impreciso, é o diâmetro da pupila: ele muda com a atividade desses neurônios, mas também com a atividade de outras regiões~\cite{Joshi2016} e das saídas colinérgicas no córtex~\cite{Reimer2016}. A pupila mostra o nível geral de excitação, e não só o \nt{NORA}.

\section{O ganho}

O que foi medido é isto: a noradrenalina suprime a atividade de fundo dos neurônios mais do que sua resposta ao estímulo, e a relação entre sinal e fundo aumenta~\cite{Foote1975NE}. Que a inclinação da característica de um neurônio isolado seja literalmente multiplicada por um coeficiente não decorre do experimento. Como no capítulo~\ref{sec:neurontypes}, tomemos um modelo simples que reproduz esse efeito: a noradrenalina muda a inclinação da característica de transferência do destinatário~\cite{ServanSchreiber1990}. As entradas fracas, sublimiares (o fundo), passam a dar ainda menos, e as fortes, ainda mais. Seja a resposta do neurônio a taxa
\begin{empheq}[box=\keybox]{equation}
r \;=\; F\bigl(g\,(u-\theta)\bigr), \qquad F(z)=\frac{r_{\max}}{1+e^{-z}} ,
\label{eq:gain}
\end{empheq}
em que $u$ é a corrente de entrada, $\theta$ é o limiar, e $g$ é o ganho, que no modelo é controlado pelo \nt{NORA}. Com $g$ pequeno, a taxa acompanha suavemente a entrada numa faixa ampla: o neurônio é um \emph{amplificador analógico}. Com $g$ grande, quase qualquer entrada acima do limiar dá a taxa máxima, e abaixo dele, zero: o neurônio é um \emph{comparador}, um elemento quase digital (fig.~\ref{fig:gaincurves}). Um único botão comuta o mesmo neurônio entre dois modos que, em eletrônica, são feitos por circuitos diferentes.

\begin{figure}[t]
\centering
\begin{tikzpicture}[>=Stealth,x=1.1cm,y=2.4cm]
\draw[->,gray!70!black] (-3.3,0) -- (3.4,0) node[below left,black] {\small $u$};
\draw[->,gray!70!black] (-3.3,0) -- (-3.3,1.2) node[left,black] {\small $r$};
\draw[densely dashed,gray!60!black] (0,0) -- (0,1.1);
\node[below,font=\small] at (0,0) {$\theta$};
\draw[densely dotted,gray!60!black] (-3.3,1) -- (3.2,1) node[right,black,font=\small] {$r_{\max}$};
\draw[very thick,blue!60!black] plot[domain=-3.2:3.2,samples=60] (\x,{1/(1+exp(-0.8*\x))});
\draw[very thick,orange!85!black] plot[domain=-3.2:3.2,samples=60] (\x,{1/(1+exp(-2*\x))});
\draw[very thick,red!70!black] plot[domain=-3.2:3.2,samples=120] (\x,{1/(1+exp(min(-8*\x,9)))});
\node[blue!60!black,font=\small,anchor=west] at (0.5,0.12) {$g$ pequeno: amplificador};
\node[red!70!black,font=\small,anchor=east] at (-0.3,0.85) {$g$ grande: comparador};
\end{tikzpicture}
\caption{Característica de transferência~\eqref{eq:gain} com três valores do ganho $g$.}
\label{fig:gaincurves}
\end{figure}

Um ganho grande ajuda contra o ruído? Nem sempre. Suponha que duas entradas diferem em $\Delta u$ e é preciso dizer qual é maior. Se o ruído $\sigma_{\text{antes}}$ é somado à entrada \emph{antes} do amplificador, amplificam-se o sinal e o ruído, e sua relação não muda. Mas se o ruído $\sigma_{\text{depois}}$ é somado \emph{depois} do amplificador (vindo de sinapses pouco confiáveis e de disparos aleatórios da própria rede, como no capítulo~\ref{sec:code}), a relação sinal-ruído
\begin{empheq}[box=\keybox]{equation}
d' \;=\; \frac{g\,\Delta u}{\sqrt{g^2\sigma_{\text{antes}}^2+\sigma_{\text{depois}}^2}}
\label{eq:gainsnr}
\end{empheq}
cresce com $g$ até que $g\sigma_{\text{antes}}$ se iguale a $\sigma_{\text{depois}}$~\cite{ServanSchreiber1990}.

\begin{keyblock}
\begin{proposition}[Quando o ganho ajuda]
\label{prop:gainnoise}
Aumentando o ganho, o \nt{NORA} melhora a discriminação das entradas só contra o ruído que surge \emph{depois} do amplificador, na própria rede. O ruído que chegou junto com a entrada, o ganho não remove.
\end{proposition}
\end{keyblock}

\section{O ganho na escolha entre duas opções}

Voltemos à escolha do capítulo~\ref{sec:gaba}: dois grupos de neurônios \nt{GLUT} com entradas $I_1$, $I_2$ se inibem mutuamente com força $\beta$. Agora cada grupo tem um ganho $g$: nas equações da escolha~\eqref{eq:wta}, a expressão entre parênteses é multiplicada por $g$. Enquanto os dois grupos funcionam, as taxas estacionárias saem de $r_1=g(I_1-\beta r_2)$, $r_2=g(I_2-\beta r_1)$. Somando e subtraindo essas igualdades, obtemos a soma $r_1+r_2=g(I_1+I_2)/(1+g\beta)$ e a diferença $r_1-r_2=g(I_1-I_2)/(1-g\beta)$. A razão entre elas diz com que nitidez a rede prefere uma entrada à outra:
\begin{empheq}[box=\keybox]{equation}
\frac{r_1-r_2}{r_1+r_2} \;=\; \varepsilon\,\frac{1+g\beta}{1-g\beta}, \qquad \varepsilon=\frac{I_1-I_2}{I_1+I_2}, \quad g\beta<1 .
\label{eq:wtagain}
\end{empheq}
Uma pequena diferença de entradas $\varepsilon$ é amplificada $(1+g\beta)/(1-g\beta)$ vezes, e esse fator cresce sem limite quando $g\beta$ se aproxima de um. Com $g\beta>1$, o estado em que os dois grupos funcionam é instável, como na proposição~\ref{prop:wta}: um vence, o outro se cala (fig.~\ref{fig:pitchfork}).

\begin{figure}[t]
\centering
\begin{tikzpicture}[>=Stealth,x=3.2cm,y=1.5cm]
\draw[->,gray!70!black] (0,0) -- (2.15,0) node[below left,black] {\small $g\beta$};
\draw[->,gray!70!black] (0,-1.25) -- (0,1.35) node[above,black] {\small $(r_1-r_2)/(r_1+r_2)$};
\draw[gray!60!black] (1,-0.04) -- (1,0.04); \node[below right,font=\small] at (1,0) {$1$};
\node[left,font=\scriptsize] at (0,1) {$1$}; \node[left,font=\scriptsize] at (0,-1) {$-1$};
\draw[very thick,blue!60!black] plot[domain=0:0.905,samples=60] (\x,{0.05*(1+\x)/(1-\x)}) -- (1,1) -- (2.05,1);
\draw[very thick,blue!60!black] (1.105,-1) -- (2.05,-1);
\draw[thick,densely dashed,gray!60!black] (1,0) -- (2.05,0);
\node[font=\small,align=center] at (0.45,-0.62) {pondera:\\os dois grupos\\funcionam};
\node[font=\small,align=center] at (1.55,0.45) {decidiu:\\só um funciona};
\node[font=\small,blue!60!black,anchor=south west] at (1.05,-0.97) {a entrada fraca venceu antes e se mantém};
\end{tikzpicture}
\caption{Escolha entre duas opções com ganhos diferentes; as entradas diferem em $\varepsilon=0.05$. Com $g\beta>1$, as duas decisões são estáveis (linhas contínuas; a vitória da entrada fraca, com $g\beta>(1+\varepsilon)/(1-\varepsilon)\approx1.1$), e a rede mantém a que já tomou; o estado simétrico (linha tracejada) é instável.}
\label{fig:pitchfork}
\end{figure}

\begin{keyblock}
\begin{proposition}[O ganho comuta o modo de escolha]
\label{prop:gainwta}
Numa rede de escolha com inibição mútua $\beta$, o ganho $g$ leva a rede a atravessar a fronteira $g\beta=1$. Abaixo dela, a rede ``pondera'': os dois grupos funcionam, e a diferença de entradas é amplificada segundo~\eqref{eq:wtagain}. Acima, a rede ``decidiu'': um grupo funciona, e a decisão se mantém. Mudando $g$, o \nt{NORA} pode comutar a rede entre esses modos sem tocar em nenhum peso.
\end{proposition}
\end{keyblock}

\section{O ganho como temperatura}

Agora acrescentemos à escolha um ruído que surge na própria rede, depois do amplificador. Qual grupo vence é decidido pelo sinal da grandeza $g(u_A-u_B)+\eta$, em que $u_A-u_B$ é a diferença de entradas, isto é, a diferença dos pesos aprendidos no capítulo~\ref{sec:dofa}, e $\eta$ é um ruído de escala $\sigma$. Se o ruído tem distribuição logística (para a gaussiana, a curva é quase a mesma), a probabilidade de escolher $A$ é
\begin{empheq}[box=\keybox]{equation}
P(A) \;=\; \frac{1}{1+e^{-g\,(u_A-u_B)/\sigma}} .
\label{eq:softmax}
\end{empheq}
O programador reconhece aqui a escolha por softmax com temperatura $T=\sigma/g$. O ganho é o inverso da temperatura. Com $g$ pequeno, mesmo uma opção claramente melhor é escolhida só um pouco mais vezes que a pior: a rede busca muito. Com $g$ grande, a melhor opção conhecida é escolhida quase sempre: a rede aproveita o que sabe.

Precisemos o que é $g$ em~\eqref{eq:wtagain} e~\eqref{eq:softmax}: é o ganho \emph{efetivo} da diferença de entradas da tarefa atual no instante da escolha. Os dois modos do \nt{NORA} agem sobre ele em sentidos opostos. A rajada curta no instante da decisão o eleva, e a rede consolida a escolha. A atividade de fundo alta, ao contrário, acompanha a busca: em humanos, antes de uma escolha ao acaso, a pupila basal é mais larga~\cite{JepmaNieuwenhuis2011}, e em ratos a entrada do \nt{NORA} no córtex cingulado anterior leva a escolha a um modo aleatório~\cite{Tervo2014stochastic}. Logo, ao \nt{NORA} de fundo alto corresponde um $g$ efetivo \emph{pequeno}, e à rajada, um $g$ grande.

\section{Quanto buscar}

Qual ganho é melhor? Verificamos isso num modelo. A rede escolhe uma de duas ações por~\eqref{eq:softmax}; cada uma traz recompensa com certa probabilidade, e a rede a estima pelas recompensas da ação escolhida, como no capítulo~\ref{sec:dofa}. De tempos em tempos o mundo muda: as duas probabilidades são sorteadas de novo. Pode mudar tanto a ação escolhida quanto aquela que a rede não experimenta há muito tempo; desta última a rede só pode saber pela busca. A fig.~\ref{fig:invertedU} mostra que fração da melhor recompensa possível a rede obtém com diferentes $g$ constantes.

\begin{figure}[t]
\centering
\begin{tikzpicture}[>=Stealth,x=1.2cm,y=1cm]
\draw[->,gray!70!black] (0,0) -- (7.6,0) node[below left,black] {\small $g$};
\draw[->,gray!70!black] (0,0) -- (0,4.4) node[above,black,font=\small] {fração da melhor recompensa};
\foreach \x/\l in {0/0.5,1/1,2/2,3/4,4/8,5/16,6/32,7/64}{\draw[gray!60!black] (\x,-0.06) -- (\x,0.06); \node[below,font=\scriptsize] at (\x,0) {\l};}
\foreach \v/\y in {0.8/0.8,0.9/2.4,1.0/4.0}{\draw[gray!60!black] (-0.06,\y) -- (0.06,\y); \node[left,font=\scriptsize] at (0,\y) {\v};}
\draw[very thick,blue!60!black] plot[mark=*,mark size=1.3pt] coordinates {(0,0.368) (1,0.880) (2,1.728) (3,2.784) (4,3.456) (5,3.616) (6,3.488) (7,3.264)};
\draw[very thick,red!70!black] plot[mark=*,mark size=1.3pt] coordinates {(0,0.368) (1,0.672) (2,1.184) (3,1.840) (4,2.176) (5,1.920) (6,1.456) (7,1.104)};
\node[blue!60!black,font=\small,anchor=south] at (5,3.7) {mundo calmo};
\node[red!70!black,font=\small,anchor=north] at (5.1,1.25) {mundo que muda depressa};
\draw[->,thick,blue!60!black] (5,3.25) -- (5,3.55);
\draw[->,thick,red!70!black] (4,1.8) -- (4,2.1);
\node[font=\small\itshape,gray!35!black,anchor=west] at (0.15,2.3) {busca às cegas};
\node[font=\small\itshape,gray!35!black] at (6.45,2.55) {não nota mudanças};
\end{tikzpicture}
\caption{Fração da melhor recompensa possível com ganho $g$ constante (escala logarítmica em $g$). Curva azul: o mundo muda em média uma vez a cada mil tentativas; vermelha, uma vez a cada vinte. As setas marcam o melhor $g$: num mundo que muda depressa, ele é duas vezes menor. Média de 60 execuções de 4000 tentativas.}
\label{fig:invertedU}
\end{figure}

Com $g=0.5$, a rede quase não aproveita o conhecimento e obtém cerca de três quartos do possível; com $g$ muito grande, perde as mudanças. O melhor valor: $g=16$ quando o mundo muda uma vez a cada mil tentativas (fração $0.976$) e $g=8$ quando muda uma vez a cada vinte ($0.886$).

\begin{keyblock}
\begin{proposition}[O U invertido]
\label{prop:explore}
A recompensa que a rede colhe depende do ganho como um U invertido: um $g$ pequeno demais gasta tentativas em busca, um grande demais não nota mudanças. O melhor $g$ é tanto menor quanto mais depressa muda o mundo.
\end{proposition}
\end{keyblock}

O U invertido também aparece nos experimentos: os animais resolvem a tarefa melhor com atividade de fundo moderada dos neurônios \nt{NORA} e respostas em rajada nítidas, e com atividade de fundo alta se distraem e alternam entre tarefas~\cite{AstonJones2005}. Isso concorda com a diferença de modos da seção anterior: a rajada no instante da decisão dá um $g$ efetivo grande exatamente quando é preciso escolher, e a atividade de fundo alta, sem rajadas, amplifica tudo, inclusive entradas alheias, de modo que o $g$ efetivo para a tarefa atual cai: é a busca.

\section{Quem gira o botão}

Já que o melhor ganho depende da rapidez com que o mundo muda, seria bom ajustá-lo. Mas como os neurônios \nt{NORA} saberiam que o mundo mudou? O indício natural é a \emph{surpresa}: os erros de predição ficaram maiores que o habitual. É importante distinguir dois tipos de incerteza. Se a recompensa é simplesmente aleatória, os erros são sempre grandes, e isso é esperado. Mas se os erros de repente cresceram em relação ao seu tamanho habitual, algo mudou. Segundo uma hipótese, é exatamente essa incerteza inesperada que o \nt{NORA} informa, e a esperada, o \nt{ACET}~\cite{YuDayan2005}.

O que fazer com esse sinal? Pode-se reduzir o ganho e buscar mais. Pode-se elevar a taxa de aprendizado, para esquecer mais depressa as estimativas obsoletas: experimentos mostram que depois de uma mudança brusca as pessoas aprendem mais depressa, e a pupila se dilata~\cite{Nassar2012}; lembremos que a pupila é indicador não só do \nt{NORA}, mas também do \nt{ACET}~\cite{Reimer2016}. E pode-se fazer as duas coisas de uma vez: segundo outra hipótese, a rajada do \nt{NORA} funciona como uma \emph{interrupção}, um sinal pelo qual a rede zera o estado atual e se reconfigura para a nova situação~\cite{BouretSara2005}, como a linha geral de interrupção de um processador.

Digamos com franqueza o que não encontramos. No modelo, testamos as duas formas simples de ajuste: reduzir $g$ ou elevar a taxa de aprendizado quando o erro suavizado passa da sua média de longo prazo. Num mundo que ora fica parado por muito tempo, ora muda com frequência, as melhores configurações das duas formas deram $0.926$--$0.927$ da melhor recompensa, quase o mesmo que o melhor $g$ constante ($0.925$ com $g=8$) ou uma taxa de aprendizado constante, mas grande o bastante ($0.928$), e as configurações ruins deram menos. As curvas da fig.~\ref{fig:invertedU} são achatadas perto do topo, e num mundo assim quase não há o que ajustar. O ajuste deve compensar onde regimes diferentes exigem configurações muito diferentes. Que lei de controle exatamente o \nt{NORA} implementa ainda não foi estabelecido.

\section{Resumo}

\begin{table}[t]
\centering
\footnotesize
\setlength{\tabcolsep}{4pt}
\renewcommand{\arraystretch}{1.15}
\begin{tabular}{>{\raggedright}p{3.0cm}>{\raggedright}p{4.4cm}>{\raggedright\arraybackslash}p{6.0cm}}
\toprule
O que o \nt{NORA} faz & Como & O que se sabe \\
\midrule
muda a inclinação da resposta dos neurônios & ganho $g$ em~\eqref{eq:gain} & aumento da relação sinal-ruído medido; o modelo multiplicativo é uma simplificação \\
comuta a escolha & leva a rede através de $g\beta=1$ & decorre do modelo~\eqref{eq:wtagain}; não há medições diretas do limiar \\
define quanto buscar & $g$ é o inverso da temperatura~\eqref{eq:softmax}; fundo: $g$ efetivo pequeno (busca); rajada: grande (decisão) & o U invertido e os dois modos de funcionamento foram medidos; a ligação com a busca é hipótese bem fundamentada \\
informa mudanças & incerteza inesperada & pupila e taxa de aprendizado crescem após mudanças: medido; que seja um sinal do \nt{NORA} é hipótese \\
interrompe e zera & rajada em resposta a evento inesperado & rajadas em eventos importantes medidas; a ``interrupção'' é hipótese \\
\bottomrule
\end{tabular}
\caption{O que o \nt{NORA} acrescenta à rede de \nt{GLUT}, \nt{GABA} e \nt{DOPA}.}
\label{tab:nora}
\end{table}

O \nt{DOPA} diz à rede \emph{para onde} mudar os pesos. O \nt{NORA} não muda nenhum peso, mas decide \emph{como} a rede usa o que já sabe (tabela~\ref{tab:nora}): um único botão, o ganho, transforma o neurônio de amplificador em comparador, a rede de escolha de ``ponderando'' em ``decidida'', e a escolha, de busca em aproveitamento. Como o \nt{NORA} descobre a rapidez com que o mundo muda é uma questão em aberto.

\section{Exercícios}

\begin{exercise}[\lvlA]
\label{ex:08:1}
\emph{Um botão para o cérebro inteiro.} (a)~Pela tabela~\ref{tab:types}, estime quantos neurônios do cérebro cabem a cada neurônio \nt{NORA}. (b)~Ruído antes do amplificador $\sigma_{\text{antes}}=0.5$, depois, $\sigma_{\text{depois}}=4$. Com que $g$ o $d'$ de~\eqref{eq:gainsnr} chega a $90\,\%$ do seu limite?
\end{exercise}

\begin{exercise}[\lvlB]
\label{ex:08:2}
\emph{Escolha com ganho.} $\tau\,\dd r_1/\dd t=-r_1+g[I_1-\beta r_2]_+$, $\tau\,\dd r_2/\dd t=-r_2+g[I_2-\beta r_1]_+$, $\tau=20\ms$. (a)~Em quanto tempo $r_1-r_2$ decai com $g\beta=0.5$ e $0.9$? (b)~Com $\varepsilon=0.05$, encontre o $g\beta$ abaixo do qual os dois grupos funcionam e acima do qual a vitória da entrada fraca é estável.
\end{exercise}

\begin{exercise}[\lvlA]
\label{ex:08:3}
\emph{Softmax a partir do ruído.} Cada um dos $K$ grupos recebe seu próprio ruído de Gumbel, $P(\eta_k<z)=\exp(-e^{-z/\sigma})$, e vence o maior $gu_k+\eta_k$. Encontre $P(k)$. Com que $g$ a melhor de duas opções com $u_A-u_B=0.1$, $\sigma=1$, é escolhida em $90\,\%$ dos casos?
\end{exercise}

\begin{exercise}[\lvlB]
\label{ex:08:4}
\emph{Estimativa do melhor ganho.} No modelo da fig.~\ref{fig:invertedU}, as probabilidades de recompensa após uma mudança são uniformes em $[0,1]$. A rede experimenta a pior ação uma vez a cada $e^{g\Delta}$ tentativas; igualando isso a $1/h$, obtemos $g^*\approx\ln(1/h)/\bar\Delta$. Encontre $\bar\Delta=\langle|p_A-p_B|\rangle$ e $g^*$ com $h=0.001$, $0.01$, $0.05$.
\end{exercise}

\begin{exercise}[\lvlB]
\label{ex:08:5}
\emph{Simulação: ganho e rapidez das mudanças.} Encontre $g^*(h)$ para $h$ de $0.001$ a $0.2$ com uma cópia de \texttt{models/\allowbreak nora\_explore.py} (ela grava arquivos na pasta atual). Onde vale a estimativa do exercício~\ref{ex:08:4}, e por que, com mudanças rápidas, $g^*$ deixa de diminuir?
\end{exercise}

\begin{exercise}[\lvlB]
\label{ex:08:6}
\emph{Projeto: interrupção para a rede de escolha.} A rede do exercício~\ref{ex:08:2} com $\beta=2$, $\tau=20\ms$, $g=1$ mantém $r_1=0.9$, $r_2=0$, embora as entradas agora sejam $I_1=0.9$, $I_2=1.0$. O \nt{NORA} reduz $g$ a $g_{\text{lo}}$ por um tempo $T_p$. Encontre o menor $T_p$ com $g_{\text{lo}}=0$ analiticamente e, com $g_{\text{lo}}=0.25$, por simulação.
\end{exercise}

\begin{exercise}[\lvlC]
\label{ex:08:7}
\emph{Quando o ajuste compensa.} Um oráculo sabe se a época é calma ou instável e usa em cada uma seu próprio $g$. (a)~Por simulação, encontre seu ganho sobre o melhor $g$ constante no mundo do capítulo (épocas de $1000$ tentativas, $h=0.001$ e $0.05$). (b)~Construa um mundo em que o ganho seja maior e encontre que parte dele o ajuste pela surpresa captura.
\end{exercise}
